\itemhead{第1章 小结与复习 · 第6题（教材p.28）}{XJ-B1-C01-SU-P028-Y1490}
\label{XJ-B1-C01-SU-P028-Y1490}
{\footnotesize 湘教版教材p.28 / PDF第35页 / 第1章 小结与复习 / 第6题}
\field{原题}{当陈述语句中涉及不确定的对象时，其真假往往难以确定。使用全称量词和存在量词可以对不确定的对象进行约束，从而使语句具有确定的真假。试结合实例讨论：如何判断含有量词的命题的真假？如何对含有一个量词的命题进行否定？}
\field{参考答案}{全称真须证明所有，否定只需一个反例；存在真只需一个实例，否定须证明全不满足。否定时交换量词、否定断言、保留论域。}
{\small \field{关联知识}{L03 全称、存在量词及真假的判定（教材帮p.38、39 / PDF47、48）；L04 含一个量词命题的否定（教材帮p.39、40 / PDF48、49）}
\field{方法与易错点}{先确认量词、作用对象及论域，再采用任意性证明或实例反例。 易错：用几个满足的数证明全称命题；把“存在”理解为“唯一存在”。}
\field{答案依据}{答案册未提供本栏目的独立对应答案；依据原题、教材概念补写参考回应。开放性总结不设唯一数值答案。}}
